\section{Escala de Dados}

Alternativamente às classificações apresentadas, podemos tipificar atributos
através de escalas de medição de seus valores. As classes são similares às
classificações apresentadas anteriormente, sendo elas:

\begin{itemize}
	\item Escala Nominal -- para dados nessa escala, só podemos dizer que seus
	      valores são diferentes de outros, sendo usada para categorizar indivíduos.
	      Um exemplo simples é o sexo de uma pessoa: feminino ou masculino. Essa
	      escala não suporta operações matemáticas, e uma medida de centralidade
	      comum para resumir dados nessa escala é a moda.

	\item Escala Ordinal -- nessa escala, os valores podem ser ordenados, e podemos
	      então dizer que, além de diferentes, um valor é maior do que o outro. Essa
	      escala tem sua estrutura preservada por operações que preservem a ordem.
	      Um exemplo de variável ordinal é o nível de escolaridade: fundamental, médio
	      e superior. Medidas de tendência central comuns para dados nessa escala são a
	      moda e a mediana.

	\item Escala Intervalar - essa escala possui uma origem arbitrária e necessita
	      de uma unidade de medida, sendo um exemplo a temperatura de um ambiente.
	      Podemos afirmar que valores são diferentes, maiores e quão maior em relação
	      a outro, e transformações afim -- do tipo $ax + b$ -- não alteram a estrutura
	      dessa escala. Podemos utilizar medidas como média, moda e mediana.

	\item Escala Razão -- nessa escala, existe uma origem absoluta, e podemos dizer que um
	      valor é o quão maior do que outro através de razões. Um exemplo de variável
	      nessa escala é o peso de um indivíduo. Transformações do tipo $ax$ preservam
	      a estrutura dessa escala, e podemos utilizar medidas como média, moda e mediana.
\end{itemize}
